\emph{Setting.} Write elements of the maximal torus of $G$ as phase tuples, in turns: four for $SU(4)$ on its $4$, two each for $SU(2)_{L}$ and $SU(2)_{R}$ on their doublets. The $(10,1,3)$ vacuum is its $\nu^{c}\nu^{c}$ component (Theorem \claimref{standard-model.majorana-element}): the weight $2e_{4}$ of the $10$ ($B - L = -2$) tensored with the weight $T_{3R} = +1$ of the $3$, so $Y = 0$ on it; a torus element fixes it exactly when the phase of that weight vanishes mod $1$. On the $4$, $B - L = \mathrm{diag}(\tfrac13, \tfrac13, \tfrac13, -1)$ (Fact \claimref{standard-model.imported-pati-salam-representation-theory-survival}).

\emph{(i)} The element $m = (-1_{4}, 1, 1)$ has phases $\tfrac12$ on the $4$; its $SU(4)$ determinant phase is $2 \equiv 0$, so $m \in G$. Every weight of the $10$ is a sum of two weights of the $4$, so $m$ acts on the $10$ by $e^{2\pi I_{4}} = 1$ (numbers and matrix entries in $\mathrm{span}\{1, I_{4}\}$, as in the statement) and trivially on the whole $(10,1,3)$; in particular $m \in H$. On the $(4,2,1)$ and $(\bar4,1,2)$ it acts by $-1$, and $\tfrac32(B - L) \equiv \tfrac12$ on every entry of the $4$, so $m = \exp(3\pi I_{4}(B - L))$ there: matter parity. In Clifford form, $SU(4) = \mathrm{Spin}(6)$ and $-1_{4}$ is the rotor $-1$, the square of the rotor $e_{1}e_{2}$ of the rotation by $\pi$; by Proposition \claimref{std-algebra.spin-double-cover} it is the kernel of $\Ad$, so it fixes every vector, hence every multivector, and the $10 \oplus \overline{10}$ is the space of $3$-vectors of the six-dimensional factor (self-dual and anti-self-dual), while on spinors (the $4$ and $\bar 4$) it acts by $-1$; the check computes these in $\Cl(6,0)$. Under $\mathrm{Spin}(6) \times \mathrm{Spin}(4) \to \mathrm{Spin}(10)$ its image is the rotor $-1$ of $\mathrm{Spin}(10)$. Of the centre $\{I_{4}^{a}1_{4}\}$ of $SU(4)$ only $a = 0, 2$ fix the vacuum (phase $a/2$ on it).

Suppose $m \in H_{0}$. A general element of $H_{0}$ is $h = (\mathrm{diag}(e^{I_{4}\theta/6}U_{3}, e^{-I_{4}\theta/2}),\, u_{L},\, \mathrm{diag}(e^{I_{4}\theta/2}, e^{-I_{4}\theta/2}))$ with $U_{3} \in SU(3)$, $u_{L} \in SU(2)_{L}$, $\theta \in \R$, since $\exp(I_{4}\theta Y)$ has these entries and $SU(3)$, $SU(2)_{L}$ and $U(1)_{Y}$ commute; $H_{0}$ is their product. Because $m$ is central in $G$, so would $h$ be. The $SU(2)_{R}$ entry $\mathrm{diag}(e^{I_{4}\theta/2}, e^{-I_{4}\theta/2})$ is a scalar only if $e^{I_{4}\theta} = 1$, that is $\theta = 2\pi t$ with $t \in \mathbb{Z}$, and then it is $(-1)^{t}$. The $SU(4)$ entry is a scalar $c$ only if $c = e^{-I_{4}\theta/2} = (-1)^{t}$ and $U_{3} = (-1)^{t}e^{-I_{4}\pi t/3}1_{3}$, a centre element of $SU(3)$ (determinant $(-1)^{3t}e^{-I_{4}\pi t} = (-1)^{4t} = 1$). So every central element of $G$ in $H_{0}$ has its $SU(4)$ scalar equal to its $SU(2)_{R}$ scalar, both $(-1)^{t}$. For $m$ these are $-1$ and $+1$: $m \notin H_{0}$. Since $G$ is connected and simply connected, the exact homotopy sequence of $H \to G \to G/H$ gives $\pi_{1}(G/H) \cong \pi_{0}(H) = H/H_{0}$ (standard, cited from memory), which contains the class of $m$, of order two ($m^{2} = 1$). The element $(-1_{4}, -1, -1)$ acts trivially on every field of the chain, so the gauge group that acts effectively is the quotient of $G$ by it; that element is central with $c_{4} = c_{R}$, hence lies in $H_{0}$, so the quotient has the same coset space $G/H$ and the same $\pi_{1}$. These are the $\mathbb{Z}_{2}$ strings; the monopoles come from $\pi_{2}(G/H) = \pi_{1}(H_{0}) \supseteq \mathbb{Z}$: the loop $\exp(2\pi I_{4}t Y)$, $t \in [0, 6]$, closes (hypercharge period $6$ on the $16$) and generates an infinite cyclic subgroup of $\pi_{1}(U(1)_{Y})$ that survives in $\pi_{1}(H_{0})$ because $U(1)_{Y}$ is a central circle factor of $H_{0}$ up to a finite quotient. A control: $(-1_{4}, 1, -1_{R})$ is the element $t = 1$, $U_{3} = e^{2\pi i/3}1_{3}$ of $H_{0}$, so $m$ and $(1, 1, -1_{R})$ represent the same nontrivial class.

\emph{(ii)} By Proposition \claimref{standard-model.what-angle-asks-neutrino-sector}, $v_{R} = M_{I}/g_{R} = 1.95 \times 10^{10}/0.469 = 4.158 \times 10^{10}$ GeV. Under Premise \claimref{neutrino.r-descent-reheats-below-lower-heavy}, with the untuned fit's cap on the lightest heavy neutrino (Proposition \claimref{synthesis.untuned-fit-caps-lightest-heavy-neutrino}), $T_{\mathrm{desc}} < 7.3 \times 10^{8}$ GeV, which lies between $0.0175\,v_{R}$ and $0.0176\,v_{R}$. A gauge symmetry broken at the scale $v_{R}$ is restored above a temperature of order $v_{R}$ (the ratio is a ratio of couplings and is not computed in the store); unless that ratio is below $0.0176$, no temperature reached after the descent restores the symmetry, so no Kibble formation of the chain's defects occurs in the descended phase, and they are present only if the descent deposits them. The remaining breakings of the series form none: $(SU(2)_{L} \times U(1)_{Y})/U(1)_{\mathrm{em}}$ is the three-sphere (cited from memory), with $\pi_{0} = \pi_{1} = \pi_{2} = 0$, and the strong interaction breaks no gauge symmetry.

\emph{(iii)} A string network of any initial density relaxes to the scaling regime, of order one string per Hubble volume, whose CMB imprint scales as $(G\mu)^{2}$ (cited from memory). For a string of the Nielsen-Olesen type $\mu = 2\pi\epsilon\,\eta^{2}$, with $\eta$ the vacuum value and $\epsilon$ a slowly varying function of the quartic-to-gauge coupling ratio, equal to one at the self-dual point and below $10$ over the range of interest (Fact \claimref{synthesis.imported-nielsen-olesen-string-tension}); with $G = 1/(8\pi M_{r}^{2})$, $G\mu = \epsilon\,\eta^{2}/(4M_{r}^{2})$. The $(10,1,3)$ vacuum value is $v_{R}$ up to the normalisation of the triplet, which is absorbed by $\eta \in [v_{R}/\sqrt2, \sqrt2\,v_{R}]$. With $(v_{R}/M_{r})^{2} = 2.9156 \times 10^{-16}$: $G\mu = 7.29 \times 10^{-17}$ at $\epsilon = 1$, $\eta = v_{R}$; $1.82 \times 10^{-17}$ at the lower corner; $1.457 \times 10^{-15}$ at the upper corner. The Planck string bounds (Fact \claimref{synthesis.imported-planck-2013-cosmic-string-bounds}) give $f_{10} < 0.015$ at $G\mu < 1.5 \times 10^{-7}$ for Nambu-Goto strings and $f_{10} < 0.028$ at $G\mu < 3.2 \times 10^{-7}$ for Abelian-Higgs strings, so $\kappa = f_{10}/(G\mu)^{2}$ is $6.7 \times 10^{11}$ and $2.7 \times 10^{11}$, both below $10^{13}$. The mirror sector, a reversal copy with the same chain and scales, at most doubles the power. Hence $f_{10} < 2 \times 10^{13} \times (1.457 \times 10^{-15})^{2} < 4.3 \times 10^{-17}$, and in units of the measured amplitude $A_{s} = 2.1 \times 10^{-9}$ (Fact \claimref{synthesis.imported-planck-2018-lcdm-parameters}) the product $f_{10}A_{s}$ is below $10^{-25}$. This is the string share of the temperature power at $\ell = 10$, an active source, not a primordial spectrum. A $\mathbb{Z}_{2}$ string is its own antistring and reconnects, so its network scales like the ones the bounds were computed for; the bound uses only the $(G\mu)^{2}$ scaling and a coefficient $\kappa$ fifteen times above the larger of the two models.

\emph{(iv)} Monopoles deposited at any temperature in the window of (ii), below $0.0176\,v_{R}$ and so well below their mass of order $4\pi v_{R}/g_{R}$, are nonrelativistic and do not scale: their density contrast is the Poisson noise of a deposit whose correlations have finite range, the class bounded by Theorem \claimref{synthesis.finite-range-descent-deposit-cannot-seed-cmb}, and their density is capped by the flux limit of Fact \claimref{synthesis.imported-macro-monopole-flux-limit}. This clause is argued.

\emph{Scope.} The result covers the topological defects of the stored chain and of its mirror copy. Global symmetries broken in the descended phase (none is stored; the chain gauges $B - L$) and sectors outside the store are not covered. If the restoration temperature of the chain were below $0.0176\,v_{R}$, strings could form after the descent, and clause (iii) would still bound their power. \emph{A link (argued):} a deposit whose Pati-Salam orientation is uncorrelated beyond the descent's Hubble length leaves a scaling $\mathbb{Z}_{2}$ network at $G\mu$ between about $2 \times 10^{-17}$ and $1.5 \times 10^{-15}$, near the reach of LISA for string backgrounds (of order $10^{-17}$, cited from memory), while an aligned deposit leaves none; this is the same long-range correlation of the deposit that a causal monopole bound on the descent history meets: a deposit whose orientation is uncorrelated beyond the descent's Hubble length leaves at least of order one monopole per Hubble volume as well.
